\documentclass[aps,preprint]{revtex4}%
\usepackage{amsfonts}
\usepackage{amsmath}
\usepackage{amssymb}
\usepackage{graphicx}%
\setcounter{MaxMatrixCols}{30}
%TCIDATA{OutputFilter=latex2.dll}
%TCIDATA{Version=5.50.0.2953}
%TCIDATA{CSTFile=revtex4.cst}
%TCIDATA{Created=Tuesday, February 19, 2008 12:58:10}
%TCIDATA{LastRevised=Tuesday, February 19, 2008 13:03:19}
%TCIDATA{<META NAME="GraphicsSave" CONTENT="32">}
%TCIDATA{<META NAME="SaveForMode" CONTENT="1">}
%TCIDATA{BibliographyScheme=Manual}
%TCIDATA{<META NAME="DocumentShell" CONTENT="Articles\SW\REVTeX 4">}
%BeginMSIPreambleData
\providecommand{\U}[1]{\protect\rule{.1in}{.1in}}
%EndMSIPreambleData
\newtheorem{theorem}{Theorem}
\newtheorem{acknowledgement}[theorem]{Acknowledgement}
\newtheorem{algorithm}[theorem]{Algorithm}
\newtheorem{axiom}[theorem]{Axiom}
\newtheorem{claim}[theorem]{Claim}
\newtheorem{conclusion}[theorem]{Conclusion}
\newtheorem{condition}[theorem]{Condition}
\newtheorem{conjecture}[theorem]{Conjecture}
\newtheorem{corollary}[theorem]{Corollary}
\newtheorem{criterion}[theorem]{Criterion}
\newtheorem{definition}[theorem]{Definition}
\newtheorem{example}[theorem]{Example}
\newtheorem{exercise}[theorem]{Exercise}
\newtheorem{lemma}[theorem]{Lemma}
\newtheorem{notation}[theorem]{Notation}
\newtheorem{problem}[theorem]{Problem}
\newtheorem{proposition}[theorem]{Proposition}
\newtheorem{remark}[theorem]{Remark}
\newtheorem{solution}[theorem]{Solution}
\newtheorem{summary}[theorem]{Summary}
\newenvironment{proof}[1][Proof]{\noindent\textbf{#1.} }{\ \rule{0.5em}{0.5em}}
\begin{document}
\preprint{TDFT II}
\title[TDFT II ]{Construction of Approximations in TDFT }
\author{First Author}
\affiliation{My Institution}
\author{Second Author}
\affiliation{My Institution}

\begin{abstract}


\end{abstract}
\startpage{1}
\maketitle
\tableofcontents


\section{Introduction}

To give the agenda of this paper requires first a succinct overview of the
preceding paper (hereafter I). It developed a constrained-search
time-dependent density functional theory (TDDFT) by (i) classification of
trace-class operators into equivalence classes labeled (in the case of density
operators) by full nuclear-electronic space-spin densities, (ii) recognition
of the set of equivalence classes as a vector bundle, and (iii) formulation of
a constructive form of constrained search both within a fiber and across the
bundle by characterization of the group properties of the contraction map
$\Xi$ (see I for definition) which generates the equivalence classes.

Step (iii) focused on pure states in $\mathcal{H}^{N}$ (notation is as in I
unless otherwise noted). The result was identification of the properties of
Lie groups of positive maps of $\mathcal{B}_{1}\left(  \mathcal{H}^{N}\right)
$ that preserve the defining relationship for the vector bundle structure,
namely that each group element transforms the kernel of $\Xi$ to itself. In
particular we described the group $\mathfrak{G}_{\ker}$,$\mathcal{\ }%
$comprised of invertible operators acting in $\mathcal{H}^{N}$, that are
diagonal in the coordinate representation, which contained the subgroups
$\mathfrak{G}_{fib}$ and $\mathfrak{G}_{den}$ of intra and interfiber
transformations respectively. The subgroup $\mathfrak{G}_{fib}$ was used to
produce a non redundant parameterization of the pure states contained in an
equivalence class $\left[  \gamma\right]  _{N}$ when acting on a reference
pure state $D_{ref}^{N}\left(  \gamma\right)  $ in $\left[  \gamma\right]
_{N}$, while the subgroup $\mathfrak{G}_{den}$ was used to give a non
redundant parameterization of pure states in the base of the vector bundle by
acting on a reference pure state $D_{ref}^{N}\left(  \gamma_{0}\right)  $. The
operators $D_{ref}^{N}\left(  \gamma\right)  $ and $D_{ref}^{N}\left(
\gamma_{0}\right)  $ are the density operators associated with the states
$|\Psi_{ref}\left(  \gamma\right)  \rangle$ and $|\Psi_{ref}\left(  \gamma
_{0}\right)  \rangle$ belonging to $\mathcal{H}^{N}$ and
\begin{align}
|\Psi_{ref}\left(  \gamma\right)  \rangle &  =e^{\mathcal{\breve{V}}%
_{den}\left(  \gamma\right)  }|\Psi_{ref}\left(  \gamma_{0}\right)
\rangle;\quad e^{\mathcal{\breve{V}}_{den}\left(  \gamma\right)  }%
\in\mathfrak{G}_{den}\nonumber\\
|\Psi\left(  \gamma\right)  \rangle &  =e^{\mathcal{\breve{V}}_{fib}%
(\gamma)+i\breve{\omega}}|\Psi_{ref}\left(  \gamma\right)  \rangle;\,\quad
e^{\mathcal{\breve{V}}_{fib}(\gamma)+i\breve{\omega}}\equiv e^{\breve
{\Upsilon}\left(  \gamma\right)  }\in\mathfrak{G}_{fib}%
\end{align}
The exponentiated integral operators are written in a shorthand notation that
was introduced in I. The subgroup $\mathfrak{G}_{fib}$ can be factorized into
the direct product $\mathcal{G}_{fib}\times\mathcal{G}_{\Xi}$ of positive,
$e^{\mathcal{\breve{V}}_{fib}}$, and $e^{i\breve{\omega}}$ unitary operators.
The coset $\mathfrak{Co}_{den}$ of the subgroup group $\mathfrak{G}_{den}$
produces an unique reference base state $|\Psi_{ref}\left(  \gamma\right)
\rangle$ from $|\Psi_{ref}\left(  \gamma_{0}\right)  \rangle$ for each fiber
$\left[  \gamma\right]  _{N}$ and the coset $\mathfrak{Co}_{fib}$ of the
subgroup $\mathfrak{G}_{fib}$ produces \emph{all }the pure states in the fiber
$\left[  \gamma\right]  _{N}$ in a unique fashion from $|\Psi_{ref}\left(
\gamma\right)  \rangle$. Given a reference state $|\Psi_{ref}\left(
\gamma_{0}\right)  \rangle$, \emph{any} pure state (in general, un-normalized)
in $\mathcal{H}^{N}$ can be parameterized in terms of these cosets as
\begin{equation}
|\Psi(\mathcal{V}_{fib}(\gamma),\omega,\mathcal{V}_{den}\left(  \gamma\right)
\rangle=e^{\mathcal{\breve{V}}_{fib}(\gamma)}\,e^{i\breve{\omega}%
}e^{\mathcal{\breve{V}}_{den}\left(  \gamma\right)  }|\Psi_{ref}\left(
\gamma_{0}\right)  \rangle\label{vv1omega1}%
\end{equation}
One can view the functions $\left\{  \mathcal{V}_{fib},\omega,\mathcal{V}%
_{den}\right\}  $ as coordinates for the manifold of pure states. One of the
most significant results of paper I (in Appendix G) was the construction of
the integral kernel $\mathcal{V}_{den}$ of the operator $\mathcal{\breve{V}%
}_{den}$ as an invertible function of the density $\gamma$ i.e. $\mathcal{V}%
_{den}=\mathcal{V}_{den}\left(  \gamma\right)  $ and thus as a transformed
coordinate of this manifold equivalent to $\gamma$ i.e.
\begin{equation}
|\Psi(\mathcal{V}_{fib}(\gamma),\omega,\mathcal{\gamma}\rangle\equiv
|\Psi(\mathcal{V}_{fib}(\gamma),\omega,\mathcal{V}_{den}\left(  \gamma\right)
\rangle
\end{equation}
details of this coordinate transformation are given in Appendix ?. Specific
forms for the non-unitary parts (in the coordinate eigenfunction
representation) and strong- orthogonality constraints on the integral kernels
of $\mathcal{\breve{V}}_{fib}(\gamma)$ and $\mathcal{\breve{V}}_{den}$ so that
they produce non rendundant parameterizations were developed in I, Appendices
F and G.

The time-dependent variational principle (TDVP) can be used to yield a set of
coupled dynamical equations equivalent to the Time Dependent Schr\"{o}dinger
Equation in terms of the state coordinates $\left\{  \mathcal{V}_{fib}%
,\omega,\mathcal{\gamma}\right\}  $ (I, eq. (54))
\begin{equation}
\{(\mathcal{V}_{fib},\omega,\gamma),\mathcal{E}\}=(\mathcal{\dot{V}}%
_{fib},\dot{\omega},\dot{\gamma}) \label{Vomegagammaeom}%
\end{equation}
This equation can be partitioned, yielding equations for Time Dependent
Density Functionals (TDDFs) , $\mathfrak{M}\left(  \mathcal{V}_{d},t\right)  $
and $\mathfrak{R}\left(  \mathcal{V}_{d},t\right)  $, of density trajectories
$\gamma\left(  t\right)  ,$ that are given in terms of trajectories $\left\{
\mathcal{V}_{fib\,\#},\omega_{\#}\right\}  $ determined by the equation
\begin{equation}
\left(
\begin{array}
[c]{cc}%
0 & \eta_{\mathcal{V}_{fib}\omega}\\
-\eta_{\mathcal{V}_{fib}\omega}^{t} & 0
\end{array}
\right)  \left(
\begin{array}
[c]{c}%
\mathcal{\dot{V}}_{fib}\\
\dot{\omega}%
\end{array}
\right)  =\left(
\begin{array}
[c]{c}%
\nabla_{\mathcal{V}_{fib}}\mathcal{E}\\
\nabla_{\omega}\mathcal{E+}\eta_{\mathcal{V}_{den}\omega}^{t}\mathcal{\dot{V}%
}_{den}%
\end{array}
\right)
\end{equation}
and an equation in terms of these functionals that determines exact density
trajectories (I, eq. (56))
\begin{equation}
\mathfrak{M}\left(  \mathcal{\gamma},t\right)  \dot{\gamma}=\mathfrak{R}%
\left(  \mathcal{V}_{d},t\right)  \label{gammadot2}%
\end{equation}
The functionals $\mathfrak{M}\left(  \mathcal{V}_{d},t\right)  $ and
$\mathfrak{R}\left(  \mathcal{V}_{d},t\right)  $ are given in I, Appendix I.

In this paper we develop approximate time dependent density functionals by
decomposing $\mathfrak{G}_{fib}$ into 1-, 2-, 3-, \ldots, particle irreducible
parts, then obtaining equations for $r^{\text{th}}$ level approximate
functionals that are determined by considering up $r$-particle parts. We
obtain these equations in terms of time dependent matrix elements of the
kinetic energy and coulomb interaction term in the respective cases. Section 3
focuses on expansions for evaluating these expressions exactly or
approximately. In Section 4 we develop and study exact equations that
determine the time- dependent density functionals, then consider
approximations to those equations and the construction of approximate
time-dependent functionals. Section 5 first treats exact equations for the
time-dependent density, with particular emphasis on the metric term introduced
by the parameterization. Consideration of approximate equations follows. We
give a brief sketch of the time-independent limit and the connection to
stationary state DFT in Section 6, then turn to Kohn-Sham type equations in
Section 7. Section 8 provides further consideration of the Born-Oppenheimer
limit and of the classical description of the nuclei. Brief summary remarks
are in Section 9.]

\section{Equations for Determining Time Dependent Density Functionsls}

\subsection{Stationarity Conditions\bigskip}

In the context of TDDFT the constrained search approach to the construction of
functionals of the density can be formulated by utilizing action functionals
defined on equivalence classes of states by
\begin{align}
\mathcal{A}_{fib}\left(  \gamma,\tau\right)   &  =\min_{D^{N}\left(  t\right)
\in\left[  \gamma\left(  t\right)  \right]  _{N}}\left\{  \mathcal{A}\left(
D^{N},\tau\right)  \right\}  \quad\nonumber\\
&  =\mathcal{A}\left(  D_{\#}^{N}\left(  \gamma\right)  ,\tau\right)
\label{restact}%
\end{align}
where $\tau=\left[  t_{initial},t_{final}\right]  $, $\left\{  \gamma
,D^{N}\right\}  $ denote the trajectories $\left\{  \gamma\left(  t\right)
,D^{N}\left(  t\right)  ;\,t\in\tau\right\}  $ with fixed initial and final
values and the full action is given by
\begin{equation}
\mathcal{A}\left(  D^{N},\tau\right)  =\int_{t_{initial}}^{t_{final}%
}\mathcal{L}\left(  D^{N}\left(  t\right)  \right)  dt
\end{equation}
in terms of the Lagrangian $\mathcal{L}\left(  D^{N}\left(  t\right)  \right)
$. The density operator trajectory, $D_{\#}^{N}\left(  t\right)  $, that
satisfies the Heisenberg Equation of Motion (HEM) for the given boundary
conditions is then given by
\begin{equation}
\mathcal{A}\left(  D_{\#}^{N}\left(  \gamma_{\#}\right)  ,\tau\right)
=\min_{\gamma\in\mathcal{P}_{N}}\left\{  \mathcal{A}\left(  D_{\#}^{N}\left(
\gamma\right)  ,\tau\right)  \right\}
\end{equation}
which also gives equations for the density trajectory $\gamma_{\#}\left(
t\right)  $ to be associated with the trajectory of states $D_{\#}^{N}\left(
\gamma_{\#}\left(  t\right)  \right)  $ i.e.
\begin{equation}
\min_{\gamma\in\mathcal{P}_{N}}\left\{  \mathcal{A}_{fib}\left(  \gamma
,\tau\right)  \right\}  =\mathcal{A}_{fib}\left(  \gamma_{\#},\tau\right)
\end{equation}
where $\mathcal{A}_{fib}\left(  \gamma,\tau\right)  =\mathcal{A}\left(
D_{\#}^{N}\left(  \gamma\right)  ,\tau\right)  $.

For pure states the full Lagrangian is
\begin{equation}
\mathcal{L}\left(  D^{N}\left(  t\right)  \right)  =\mathcal{K}\left(
D^{N}\left(  t\right)  \right)  -\mathcal{E}\left(  D^{N}\left(  t\right)
\right)  \label{LA}%
\end{equation}
where
\begin{equation}
\mathcal{K}\left(  D^{N}\left(  t\right)  \right)  =\frac{i}{2}\frac{\left(
\left\langle \Psi\left(  t\right)  |\dot{\Psi}\left(  t\right)  \right\rangle
-\left\langle \dot{\Psi}\left(  t\right)  |\Psi\left(  t\right)  \right\rangle
\right)  }{\left\langle \Psi\left(  t\right)  |\Psi\left(  t\right)
\right\rangle }\equiv\mathcal{K}\left(  \gamma\left(  t\right)  ,\mathcal{V}%
_{fib}\left(  t\right)  ,\omega\left(  t\right)  \right)
\end{equation}
and
\begin{equation}
\mathcal{E}\left(  D^{N}\left(  t\right)  \right)  =\frac{\left\langle
\Psi\left(  t\right)  |H\Psi\left(  t\right)  \right\rangle }{\left\langle
\Psi\left(  t\right)  |\Psi\left(  t\right)  \right\rangle }\equiv
\mathcal{E}\left(  \gamma\left(  t\right)  ,\mathcal{V}_{fib}\left(  t\right)
,\omega\left(  t\right)  \right)
\end{equation}
If we limit attention to pure states $D^{N}\left(  t\right)  =\left|
\Psi\left(  t\right)  \right\rangle \left\langle \Psi\left(  t\right)
\right|  $ the action associated with each fiber becomes
\begin{equation}
\mathcal{A}_{fib}\left(  \gamma,\tau\right)  =\mathcal{A}\left(
\gamma,\mathcal{V}_{fib}^{\#},\mathcal{V}_{fib}^{\#},\tau\right)
=\int_{t_{initial}}^{t_{final}}\mathcal{L}_{fib}\left(  \gamma\left(
t\right)  ,\mathcal{V}_{fib}^{\#}\left(  t\right)  ,\omega^{\#}\left(
t\right)  \right)  dt \label{fact}%
\end{equation}
where the Lagrangian, $\mathcal{L}_{fib}$, is given by eq. $\left(
\ref{LA}\right)  $ where now
\begin{equation}
\Psi\left(  t\right)  =\Psi\left(  \gamma\left(  t\right)  ,\mathcal{V}%
_{fib}^{\#}\left(  t\right)  ,\omega^{\#}\left(  t\right)  \right)
\end{equation}
which can be expanded to give
\begin{align}
\mathcal{L}_{fib}\left(  \gamma,\mathcal{V}_{fib}^{\#},\omega^{\#}\right)   &
=\int\mathcal{\dot{\gamma}}\left(  \mathbf{z}\right)  \mathcal{Z}\left(
\mathcal{\gamma}\right)  \left(  \mathbf{z}\right)  d\mathbf{z+}\nonumber\\
&  \int\left\{  \mathcal{\dot{V}}_{fib}^{\#}\left(  \mathsf{z}\right)
\mathcal{Z}\left(  \mathcal{V}_{fib}^{\#}\right)  \left(  \mathsf{z}\right)
+\mathcal{\dot{\omega}}^{\#}\left(  \mathsf{z}\right)  \mathcal{Z}\left(
\mathcal{\omega}^{\#}\right)  \left(  \mathsf{z}\right)  \right\}
d\mathsf{z}-\mathcal{E}\left(  \gamma,\mathcal{V}_{fib}^{\#},\mathcal{\omega
}^{\#}\right)
\end{align}
at each point of time $t$, (recall $\mathsf{z}=\kappa_{1},\ldots,\kappa
_{N}=\mathbf{x}_{1},\ldots,\mathbf{x}_{N_{e}},\mathbf{X}_{1},\ldots
,\mathbf{X}_{N_{n}}$ and $\mathbf{z}\equiv(\mathbf{x}\,,\mathbf{X})$), where
\begin{equation}
\mathcal{Z}\left(  \mathcal{\gamma}\right)  \left(  \mathbf{z}\right)
=\frac{i}{2}\left\langle \Psi|\Psi\right\rangle ^{-1}\left\{  \left\langle
\Psi\left|  \frac{\delta\Psi}{\delta\mathcal{\gamma}\left(  \mathbf{z}\right)
}\right.  \right\rangle -\left\langle \left.  \frac{\delta\Psi}{\delta
\mathcal{\gamma}\left(  \mathbf{z}\right)  }\right|  \Psi\right\rangle
\right\}
\end{equation}
and the derrivatives wrt $\gamma\left(  \mathbf{z}\right)  $ are given by the
coordinate transformation $\frac{\delta}{\delta\gamma\left(  \mathbf{z}%
\right)  }=\int\frac{\delta\mathcal{V}_{den}\left(  \mathsf{z}\right)
}{\delta\gamma\left(  \mathbf{z}\right)  }\frac{\delta}{\delta\mathcal{V}%
_{den}\left(  \mathsf{z}\right)  }d\mathsf{z}$%
\begin{align}
\mathcal{Z}\left(  \mathcal{V}_{fib}^{\#}\right)  \left(  \mathsf{z}\right)
&  =\frac{i}{2}\left\langle \Psi|\Psi\right\rangle ^{-1}\left\{  \left\langle
\Psi\left|  \frac{\delta\Psi}{\delta\mathcal{V}_{fib}\left(  \mathsf{z}%
\right)  }\right.  \right\rangle -\left\langle \left.  \frac{\delta\Psi
}{\delta\mathcal{V}_{fib}\left(  \mathsf{z}\right)  }\right|  \Psi
\right\rangle \right\} \nonumber\\
\mathcal{Z}\left(  \omega^{\#}\right)  \left(  \mathsf{z}\right)   &
=\frac{i}{2}\left\langle \Psi|\Psi\right\rangle ^{-1}\left\{  \left\langle
\Psi\left|  \frac{\delta\Psi}{\delta\omega\left(  \mathsf{z}\right)  }\right.
\right\rangle -\left\langle \left.  \frac{\delta\Psi}{\delta\omega\left(
\mathsf{z}\right)  }\right|  \Psi\right\rangle \right\}
\end{align}
and
\begin{equation}
\left|  \Psi\right\rangle \equiv\left|  \Psi\left(  \gamma\left(  t\right)
,\mathcal{V}_{fib}^{\#}\left(  t\right)  ,\omega^{\#}\left(  t\right)
\right)  \right\rangle
\end{equation}
We show below that
\begin{equation}
\mathcal{Z}\left(  \mathcal{\gamma}\right)  \left(  \mathbf{z}\right)
=0\text{ and }\mathcal{Z}\left(  \mathcal{V}_{fib}\right)  \left(
\mathsf{z}\right)  =0
\end{equation}
so that the Lagrangian $\mathcal{L}_{fib}$ becomes
\begin{equation}
\mathcal{L}_{fib}\left(  \gamma,\mathcal{V}_{fib}^{\#},\omega^{\#}\right)
=\int\mathcal{\dot{\omega}}^{\#}\left(  \mathsf{z}\right)  \mathcal{Z}\left(
\mathcal{\omega}^{\#}\right)  \left(  \mathsf{z}\right)  d\mathsf{z}%
-\mathcal{E}\left(  \gamma,\mathcal{V}_{fib}^{\#},\mathcal{\omega}%
^{\#}\right)
\end{equation}
This Lagrangian together with the stationary action condition $\delta
\mathcal{A}_{fib}=0$ produces the equation of motion eq. $\left(
\ref{gammadot2}\right)  $for the exact density trajectory $\gamma^{\#}\left(
t\right)  $.

The constrained Action and Lagrangian used in eq. $\left(  \ref{restact}%
\right)  $ to obtain the fiber Action and Lagrangian are
\begin{align}
\mathcal{L}\left|  _{fib}\left(  \gamma\left(  t\right)  ,\mathcal{V}%
_{fib}\left(  t\right)  ,\omega\left(  t\right)  \right)  \right.   &
=\nonumber\\
\mathcal{A}\left|  _{fib}\left(  \gamma,\mathcal{V}_{fib},\omega,\tau\right)
\right.   &  =\int_{t_{initial}}^{t_{final}}\mathcal{L}\left|  _{fib}\left(
\gamma\left(  t\right)  ,\mathcal{V}_{fib}\left(  t\right)  ,\omega\left(
t\right)  \right)  \right.  dt
\end{align}
together with the constraint
\begin{equation}
\end{equation}
and the fiber Action is given by
\begin{equation}
\delta\mathcal{A}\left|  _{fib}\left(  \gamma,\mathcal{V}_{fib},\omega
,\tau\right)  \right.  =0
\end{equation}
We discuss the form of these Lagrangians in more detail in paper III of this
series within the context of path integrals.

The action principle in eq. $\left(  \ref{restact}\right)  $ leads to Euler -
Lagrange equations of motion that determine trajectories of states $\left\{
D_{\#}^{N}\left(  \gamma\right)  \right\}  $ for each density $\gamma\left(
t\right)  $, however the solutions $D_{\#}^{N}\left(  \gamma\right)  $ are
only determined up to positive time dependent (no space or spin dependence)
multiplicative factors $\kappa\left(  t\right)  $, i.e. $\kappa\left(
t\right)  D_{\#}^{N}\left(  \gamma\left(  t\right)  \right)  $ is also a
solution as the factor $\kappa\left(  t\right)  $ just adds a total time
derrivative to Lagrangian $\mathcal{L}\left(  D^{N}\left(  t\right)  \right)
$. The solution trajecories $\left\{  D_{\#}^{N}\left(  \gamma\right)
\right\}  $ determine time dependent density functionals $F_{\mathcal{K}}$ and
$F_{\mathcal{E}} $ of time dependent densities
\begin{align}
F_{\mathcal{K}}\left(  \gamma\left(  t\right)  \right)   &  =\mathcal{K}%
\left(  D_{\#}^{N}\left(  \gamma\left(  t\right)  \right)  \right) \nonumber\\
F_{\mathcal{E}}\left(  \gamma\left(  t\right)  \right)   &  =\mathcal{E}%
\left(  D_{\#}^{N}\left(  \gamma\left(  t\right)  \right)  \right)
\end{align}
The functional $F_{\mathcal{E}}\left(  \gamma\left(  t\right)  \right)  $ does
not depend on $\kappa\left(  t\right)  $, while one must make a fixed choice
for $F_{\mathcal{K}}\left(  \gamma\left(  t\right)  \right)  $.

The functional $F_{\mathcal{E}}$ can be decomposed into electron and nuclear
kinetic energy, electron and nuclear coulomb interaction and electron-nuclear
coulomb interaction. The nuclear and electron coulomb interaction terms can be
further decomposed into classical coulomb and exchange-correlation parts. We
emphasize that there is \emph{no }external potential energy term in the
Hamiltonian $H$.

The condition for eq. $\left(  \ref{fibact}\right)  $ to be satisfied is that
the first order restricted variation vanish i.e.
\begin{equation}
\delta\mathcal{L}\left(  \Upsilon\left(  \gamma\left(  t\right)  \right)
,\gamma\left(  t\right)  \right)  =0;\quad\text{fixed trajectory }%
\gamma\label{var1}%
\end{equation}
and that the second order restricted variation be positive i.e.
\begin{equation}
\delta^{2}\mathcal{L}\left(  \Upsilon\left(  \gamma\left(  t\right)  \right)
,\gamma\left(  t\right)  \right)  \geq0;\quad\text{fixed trajectory }\gamma
\end{equation}
The constrained stationarity condition eq. $\left(  \ref{var1}\right)  $ for a
fixed $\left\{  \gamma\left(  t\right)  ;\,t\in\tau\right\}  $ leads to the
equations of motion described in Appendix I of paper I
\begin{equation}
\left(
\begin{array}
[c]{c}%
\mathcal{\dot{V}}_{fib}\\
\dot{\omega}%
\end{array}
\right)  =\eta\left(  \mathcal{\gamma}\right)  ^{-1}\left(
\begin{array}
[c]{c}%
\nabla_{\mathcal{V}_{fib}}\mathcal{E}\\
\nabla_{\omega}\mathcal{E}+\eta_{\mathcal{V}_{den}\omega}^{t}\mathcal{\dot{V}%
}_{den}%
\end{array}
\right)  +\sum_{1\leq i\leq m_{1}}\lambda_{i}\left(  \mathcal{\gamma}\right)
\left|  k_{i}\left(  \mathcal{\gamma}\right)  \right\rangle \label{EOM}%
\end{equation}
and
\begin{equation}
\eta\left(  \mathcal{\gamma}\right)  =\left(
\begin{array}
[c]{cc}%
0 & \eta_{\mathcal{V}_{fib}\omega}\\
\eta_{\omega\mathcal{V}_{fib}} & 0
\end{array}
\right)  =\left(
\begin{array}
[c]{ccc}%
0 & \eta_{\mathcal{V}_{fib}\omega_{den}} & \eta_{\mathcal{V}_{fib}\omega
_{fib}}\\
\eta_{\omega_{den}\mathcal{V}_{fib}} & 0 & 0\\
\eta_{\omega_{fib}\mathcal{V}_{fib}} & 0 & 0
\end{array}
\right)  \label{etamat}%
\end{equation}
for $A=\mathcal{\Upsilon}^{\ast}$, $\mathcal{\Upsilon}$ or $\gamma$;
$B=\mathcal{\Upsilon}^{\ast}$, $\mathcal{\Upsilon}$ or $\gamma$ and
$\mathtt{z}=\mathsf{z}$ in the cases $\mathcal{\Upsilon}^{\ast}$and
$\mathcal{\Upsilon}$ while \texttt{z}$=\mathbf{z}$ in the case of $\gamma$,
and the overlap $S$ is

In terms of these variables we obtain four types of derrivatives, namely%
\begin{align}
&  i\left\{  \frac{\delta}{\delta\mathcal{\tilde{V}}_{fib}(\mathsf{z})}%
\frac{\delta}{\delta\mathcal{V}_{fib}(\mathsf{z}^{\prime})}\,-\,\frac{\delta
}{\delta\mathcal{V}_{fib}(\mathsf{z})}\frac{\delta}{\delta\mathcal{\tilde{V}%
}_{fib}(\mathsf{z}^{\prime})}\right\}  \ln S|_{\mathcal{\tilde{V}}%
_{fib}=\mathcal{V}_{fib}}\nonumber\\
&  i\left\{  \frac{\delta}{\delta\mathcal{\tilde{\omega}}(\mathsf{z})}%
\frac{\delta}{\delta\mathcal{\omega}(\mathsf{z}^{\prime})}\,-\,\frac{\delta
}{\delta\mathcal{\omega}(\mathsf{z})}\frac{\delta}{\delta\mathcal{\tilde
{\omega}}(\mathsf{z}^{\prime})}\right\}  \ln S|_{\mathcal{\tilde{\omega}%
}=\mathcal{\omega}} \label{same}%
\end{align}
and%
\begin{align}
&  i\left\{  \frac{\delta}{\delta\mathcal{\tilde{V}}_{fib}(\mathsf{z})}%
\frac{\delta}{\delta\mathcal{\omega}(\mathsf{z}^{\prime})}\,-\,\frac{\delta
}{\delta\mathcal{V}_{fib}(\mathsf{z})}\frac{\delta}{\delta\mathcal{\tilde
{\omega}}(\mathsf{z}^{\prime})}\right\}  \ln S|_{\mathcal{\tilde{V}}%
_{fib}=\mathcal{V}_{fib},\mathcal{\tilde{\omega}}=\mathcal{\omega}}\nonumber\\
&  i\left\{  \frac{\delta}{\delta\mathcal{\tilde{\omega}}(\mathsf{z})}%
\frac{\delta}{\delta\mathcal{V}_{fib}(\mathsf{z}^{\prime})}\,-\,\frac{\delta
}{\delta\mathcal{\omega}(\mathsf{z})}\frac{\delta}{\delta\mathcal{\tilde{V}%
}_{fib}(\mathsf{z}^{\prime})}\right\}  \ln S|_{\mathcal{\tilde{V}}%
_{fib}=\mathcal{V}_{fib},\mathcal{\tilde{\omega}}=\mathcal{\omega}}
\label{mixed}%
\end{align}
we show shortly that derrivatives eq. $\left(  \ref{same}\right)  $ are zero
while the ones in eq. $\left(  \ref{mixed}\right)  $ are skew symmetric with
respect to each other. Hence in terms of the real and imaginary parts
\begin{equation}
\eta_{\mathbf{\digamma\,\digamma}}(t)=\left(
\begin{array}
[c]{cc}%
0 & \eta_{\mathcal{V}_{fib}\omega}\left(  t\right) \\
\eta_{\omega\mathcal{V}_{fib}}\left(  t\right)  & 0
\end{array}
\right)  =\left(
\begin{array}
[c]{cc}%
0 & \eta_{\mathcal{V}_{fib}\omega}\left(  t\right) \\
-\eta_{\mathcal{V}_{fib}\omega}\left(  t\right)  ^{t} & 0
\end{array}
\right)  \label{reim}%
\end{equation}
where $\eta_{\mathcal{V}_{fib}\omega}\left(  t,\mathsf{z},\mathsf{z}^{\prime
}\right)  =-\eta_{\omega\mathcal{V}_{fib}}\left(  t,\mathsf{z}^{\prime
},\mathsf{z}\right)  $ and $\eta_{\mathcal{V}_{fib}\omega}\left(
t,\mathsf{z},\mathsf{z}^{\prime}\right)  ^{t}=\eta_{\mathcal{V}_{fib}\omega
}\left(  t,\mathsf{z}^{\prime},\mathsf{z}\right)  $. All the quantities in eq.
$\left(  \ref{var1}\right)  $ to eq. $\left(  \ref{reim}\right)  $ refer to
the trajectory $\gamma$ at a specific time $t$. Taking into account the form
of $\eta_{\mathbf{\digamma\,\digamma}}(t)$ given in eq. $\left(
\ref{reim}\right)  $ eq. $\left(  \ref{EOM}\right)  $ can be expressed as
\begin{align}
\eta_{\mathcal{V}_{fib}\omega}\dot{\omega}  &  =\nabla_{\mathcal{V}_{fib}%
}\mathcal{E}\left(  t\right) \\
\eta_{\mathcal{V}_{fib}\omega}^{t}\mathcal{\dot{V}}_{fib}  &  =-\nabla
_{\mathcal{\omega}}\mathcal{E}\left(  t\right)
\end{align}
The trajectory of parameters $\left(  \mathcal{V}_{fib\#}(\gamma\left(
t\right)  ),\omega_{\#}\left(  t\right)  \right)  $ controls which pure state
in the fiber $\left[  \gamma\left(  t\right)  \right]  _{N}$ is selected by
the EOM eq. $\left(  \ref{EOM}\right)  $ and thus determines the time
dependent functionals.

The first step is to determine the trajectory $\left(  \mathcal{V}%
_{fib\#}\left(  \gamma\right)  ,\omega_{\#}\left(  \gamma\right)  \right)  $
for a specific Hamiltonian $H$ and thus obtain time depnendent density
functionals. The Hamiltonian $H$ contains parts that are common to all
electron-nuclei systems of a given number of electrons, namely the electron
kinetic and coulomb interaction terms, and this part defines a true universal
functional of the density. However the nuclear part and nuclear-electron
interaction is clearly system dependent as the mass and charge of nuclei vary.
The equation that determines $\left(  \mathcal{V}_{fib\#}\left(
\gamma\right)  ,\omega_{\#}\left(  \gamma\right)  \right)  $

\begin{itemize}
\item contains the metric terms

$\eta_{\mathcal{V}_{fib}\omega_{fib}}$,$\eta_{\mathcal{\omega}_{fib}%
\mathcal{V}_{den}}$,$\eta_{\omega_{fib}\mathcal{V}_{fib}}$,$\eta
_{\mathcal{V}_{fib}\mathcal{\omega}_{den}}$,$\eta_{\omega_{den}\mathcal{V}%
_{den}}$,$\eta_{\mathcal{V}_{den}\omega_{fib}}$,$\eta_{\mathcal{V}_{fib}%
\omega_{den}}$

\item the force terms $\nabla_{\omega_{fib}}\mathcal{E}$,$\nabla
_{\mathcal{V}_{fib}}\mathcal{E}$,$\nabla_{\mathcal{V}_{den}}\mathcal{E}$ and

\item the \textquotedblright velocity terms\textquotedblright\ $\mathcal{\dot
{V}}_{den}$,$\dot{\omega}_{den}$,$\mathcal{\dot{V}}_{fib}$,$\dot{\omega}%
_{fib}$,$\dot{\omega}_{den}$
\end{itemize}

all of which are evaluated along a fixed density trajectory $\left\{
\gamma\left(  t\right)  ;\,t\in\tau\right\}  $

\subsection{\bigskip Metric Terms}

The metric terms all have the basic form
\begin{equation}
\eta_{\lambda_{1},\lambda_{2}}(t,\mathsf{z},\mathsf{z}^{\prime})=i\left\{
\frac{\delta}{\delta\tilde{\lambda}_{1}(\mathsf{z})}\frac{\delta}%
{\delta\lambda_{2}(\mathsf{z}^{\prime})}-\frac{\delta}{\delta\lambda
_{2}(\mathsf{z})}\frac{\delta}{\delta\tilde{\lambda}_{1}(\mathsf{z}^{\prime}%
)}\right\}  \ln S| _{\substack{\tilde{\lambda}_{1}=\lambda_{1} \\\tilde
{\lambda}_{2}=\lambda_{2} }} \label{metric}%
\end{equation}
where $\lambda_{1},\lambda_{2}\in\left\{  \mathcal{V}_{fib},\mathcal{V}%
_{den},\omega_{fib},\omega_{den}\right\}  $. The overlap term $S$ is given by
\begin{align}
S  &  =\left\langle \widetilde{e^{\mathcal{\breve{V}}_{fib}(\mathcal{V}%
_{den})}e^{i\mathcal{\breve{\omega}}_{fib}(\mathcal{V}_{den})}%
e^{i\mathcal{\breve{\omega}}_{den}(\mathcal{V}_{den})}e^{\mathcal{\breve{V}%
}_{den}}}\Psi_{ref}\left(  \gamma_{0}\right)  \right. \nonumber\\
&  \qquad\qquad\qquad\left.  |e^{\mathcal{\breve{V}}_{fib}(\mathcal{V}_{den}%
)}e^{i\mathcal{\breve{\omega}}_{fib}(\mathcal{V}_{den})}e^{i\mathcal{\breve
{\omega}}_{den}(\mathcal{V}_{den})}e^{\mathcal{\breve{V}}_{den}}\Psi
_{ref}\left(  \gamma_{0}\right)  \right\rangle
\end{align}
where $\mathcal{V}_{den}$ is a fixed trajectory that corresponds to a density
trajectory $\gamma\left(  t\right)  $. Without loss of generality we can take
the reference state to be an IPS
\begin{equation}
\left|  \Psi_{ref}\left(  \gamma_{0}\right)  \right\rangle =\left|  \psi
_{1}\text{\textperiodcentered\textperiodcentered}\psi_{N_{e}}\right\rangle
\otimes\left|  \varphi_{1}\text{\textperiodcentered\textperiodcentered}%
\varphi_{N_{n}}\right\rangle
\end{equation}
and choose $e^{i\mathcal{\breve{\omega}}_{den}(\mathcal{V}\left(  t\right)
_{den})}e^{\mathcal{\breve{V}}\left(  t\right)  _{den}}\left|  \Psi
_{ref}\left(  \gamma_{0}\left(  t\right)  \right)  \right\rangle $ to be an
IPS\ for all times $t\in\tau$. It is not necessary to make this choice,
trajectories of reference states can also be chosen to belong other families
of pure states. However in general other choices will result in more
complicated expectation values, though could produce better low order
approximations to the time dependent density functionals. Recall that
\begin{align}
&  \mathcal{T}\left(  t\right)  \mathcal{=}\int\left(  \mathcal{V}%
_{fib}\left(  \mathcal{V}_{den}\left(  t\right)  \right)  (\mathbf{x}%
_{1},..,\mathbf{x}_{N_{e}}\mathbf{X}_{1},..,\mathbf{X}_{N_{n}})+\right.
\nonumber\\
&  \mathcal{V}_{den}\left(  t,\mathbf{x}_{1},\ldots,\mathbf{x}_{N_{e}%
}\mathbf{X}_{1},..,\mathbf{X}_{N_{n}}\right)  \left.  +i\mathcal{\omega
}\left(  \mathcal{V}_{den}\left(  t\right)  \right)  (\mathbf{x}%
_{1},..,\mathbf{x}_{N_{e}}\mathbf{X}_{1},..,\mathbf{X}_{N_{n}})\right)
\nonumber\\
&  \qquad\qquad\qquad\times\prod\limits_{j=1}^{N_{e}}\Phi_{e}^{\dagger}%
(x_{j})\Phi_{e}(x_{j})\prod\limits_{k=1}^{N_{n}}\Phi_{n}^{\dagger}(X_{k}%
)\Phi_{n}(X_{k})\prod\limits_{l=1}^{N_{e}}dx_{l}\prod_{m=1}^{N_{n}}dX_{m}%
\end{align}
where $\mathcal{V}_{fib}\left(  \mathcal{V}_{den}\left(  t\right)  \right)
=\frac{1}{2}\ln\left(  1+\phi\left(  \mathcal{V}_{den}\left(  t\right)
\right)  \right)  $ and $\phi\in L_{\infty}\left(  \mathtt{Z}\right)  $,
$\perp_{s}K_{D}$ and biorthogonal to the integral kernels of all super
operators belonging to $S_{D}$.

The second derrivatives in eq. $\left(  \ref{metric}\right)  $ have the same
form
\begin{equation}
\frac{\delta}{\delta\tilde{\lambda}_{1}(\mathsf{z})}\frac{\delta}%
{\delta\lambda_{2}(\mathsf{z}^{\prime})}\ln S=\frac{1}{S}\left\{  \frac
{\delta^{2}S}{\delta\tilde{\lambda}_{1}(\mathsf{z})\delta\lambda
_{2}(\mathsf{z}^{\prime})}-\frac{1}{S}\frac{\delta S}{\delta\tilde{\lambda
}_{1}(\mathsf{z})}\frac{\delta S}{\delta\lambda_{2}(\mathsf{z}^{\prime}%
)}\right\}
\end{equation}
different combinations of $\lambda_{1},\lambda_{2}\in\left\{  \mathcal{V}%
_{fib},\mathcal{V}_{den},\omega_{fib},\omega_{den}\right\}  $ result in
\begin{align}
&  \frac{\delta}{\delta\tilde{\lambda}_{1}(\mathsf{z})}\frac{\delta}%
{\delta\lambda_{2}(\mathsf{z}^{\prime})}\ln S=\frac{1}{S}\left\{  \left\langle
\Psi_{ref}\left(  \gamma_{0}\right)  |\Lambda\left(  \mathsf{z}\right)
\Lambda\left(  \mathsf{z}^{\prime}\right)  \Psi_{ref}\left(  \gamma
_{0}\right)  \right\rangle \right. \nonumber\\
&  \qquad\qquad\left.  \left\langle \Psi_{ref}\left(  \gamma_{0}\right)
|\Lambda\left(  \mathsf{z}\right)  \Psi_{ref}\left(  \gamma_{0}\right)
\right\rangle \left\langle \Psi_{ref}\left(  \gamma_{0}\right)  |\Lambda
\left(  \mathsf{z}^{\prime}\right)  \Psi_{ref}\left(  \gamma_{0}\right)
\right\rangle \right\}
\end{align}
where%
\begin{equation}
\Lambda\left(  \mathsf{z}\right)  =\prod\limits_{j=1}^{N_{e}}\Phi_{e}%
^{\dagger}(x_{j})\Phi_{e}(x_{j})\prod\limits_{k=1}^{N_{n}}\Phi_{n}^{\dagger
}(X_{k})\Phi_{n}(X_{k})
\end{equation}
Letting leading to the relationships $\mathcal{V=V}_{den}$ or $\mathcal{V}%
_{fib}$ and $\omega=\omega_{den}$ or $\omega_{fib\text{ }}$ one obtains
\begin{align}
\frac{\delta}{\delta\mathcal{\tilde{V}}(\mathsf{z})}\frac{\delta}{\delta
\omega(\mathsf{z}^{\prime})}\ln S  &  =i\frac{\delta}{\delta\mathcal{\tilde
{V}}(\mathsf{z})}\frac{\delta}{\delta\mathcal{V}(\mathsf{z}^{\prime})}\ln S\\
\frac{\delta}{\delta\mathcal{\tilde{\omega}}(\mathsf{z})}\frac{\delta}%
{\delta\mathcal{V}(\mathsf{z}^{\prime})}\ln S  &  =-i\frac{\delta}%
{\delta\mathcal{\tilde{V}}(\mathsf{z})}\frac{\delta}{\delta\mathcal{V}%
(\mathsf{z}^{\prime})}\ln S\\
\frac{\delta}{\delta\mathcal{\tilde{\omega}}(\mathsf{z})}\frac{\delta}%
{\delta\omega(\mathsf{z}^{\prime})}\ln S  &  =i\frac{\delta}{\delta
\mathcal{\tilde{V}}(\mathsf{z})}\frac{\delta}{\delta\mathcal{V}(\mathsf{z}%
^{\prime})}\ln S
\end{align}
Hence one can see that
\begin{align}
\frac{\delta}{\delta\mathcal{\tilde{V}}(\mathsf{z})}\frac{\delta}%
{\delta\mathcal{V}(\mathsf{z}^{\prime})}\ln S  &  =\frac{\delta}%
{\delta\mathcal{V}(\mathsf{z})}\frac{\delta}{\delta\mathcal{\tilde{V}%
}(\mathsf{z}^{\prime})}\ln S\\
\frac{\delta}{\delta\mathcal{\tilde{V}}(\mathsf{z})}\frac{\delta}%
{\delta\mathcal{\omega}(\mathsf{z}^{\prime})}\ln S  &  =-\frac{\delta}%
{\delta\mathcal{V}(\mathsf{z})}\frac{\delta}{\delta\mathcal{\tilde{\omega}%
}(\mathsf{z}^{\prime})}\ln S\\
\frac{\delta}{\delta\mathcal{\tilde{\omega}}(\mathsf{z})}\frac{\delta}%
{\delta\mathcal{V}(\mathsf{z}^{\prime})}\ln S  &  =-\frac{\delta}{\delta
\omega(\mathsf{z})}\frac{\delta}{\delta\mathcal{\tilde{V}}(\mathsf{z}^{\prime
})}\ln S\\
\frac{\delta}{\delta\mathcal{\tilde{\omega}}(\mathsf{z})}\frac{\delta}%
{\delta\omega(\mathsf{z}^{\prime})}\ln S  &  =i\frac{\delta}{\delta
\omega(\mathsf{z})}\frac{\delta}{\delta\mathcal{\tilde{\omega}}(\mathsf{z}%
^{\prime})}\ln S
\end{align}
note that in all cases the order in which the functional derrivatives are
taken is immaterial i.e.%
\begin{equation}
\frac{\delta}{\delta\tilde{\lambda}_{1}(\mathsf{z})}\frac{\delta}%
{\delta\lambda_{2}(\mathsf{z}^{\prime})}\ln S=\frac{\delta}{\delta\lambda
_{2}(\mathsf{z}^{\prime})}\frac{\delta}{\delta\tilde{\lambda}_{1}(\mathsf{z}%
)}\ln S
\end{equation}
as one would expect from the continuous function $\ln S.$Thus the various
components of the metric term are
\begin{align}
\eta_{\mathcal{V}_{fib}\mathcal{V}_{fib}}  &  =\eta_{\mathcal{V}%
_{fib}\mathcal{V}_{den}}=\eta_{\mathcal{V}_{den}\mathcal{V}_{fib}}%
=\eta_{\mathcal{V}_{den}\mathcal{V}_{den}}=0\\
\eta_{\mathcal{\omega}_{fib}\mathcal{\omega}_{fib}}  &  =\eta_{\mathcal{\omega
}_{fib}\mathcal{\omega}_{den}}=\eta_{\mathcal{\omega}_{den}\mathcal{\omega
}_{fib}}=\eta_{\mathcal{\omega}_{den}\mathcal{\omega}_{den}}=0
\end{align}%
\begin{align}
\eta_{\mathcal{V\omega}}  &  =\frac{2i}{\mathcal{N}\left(  \Upsilon\left(
\gamma\left(  t\right)  \right)  ,\gamma\left(  t\right)  \right)  }\left\{
\left\langle \Psi_{ref}\left(  \gamma_{0}\right)  |\Lambda\left(
\mathsf{z}\right)  \Lambda\left(  \mathsf{z}^{\prime}\right)  \Psi
_{ref}\left(  \gamma_{0}\right)  \right\rangle \right. \nonumber\\
&  \qquad\qquad\left.  \left\langle \Psi_{ref}\left(  \gamma_{0}\right)
|\Lambda\left(  \mathsf{z}\right)  \Psi_{ref}\left(  \gamma_{0}\right)
\right\rangle \left\langle \Psi_{ref}\left(  \gamma_{0}\right)  |\Lambda
\left(  \mathsf{z}^{\prime}\right)  \Psi_{ref}\left(  \gamma_{0}\right)
\right\rangle \right\} \nonumber\\
&  =-\eta_{\omega\mathcal{V}} \label{vomega}%
\end{align}
One can thus write $\eta_{\mathcal{V\omega}}$ as
\begin{equation}
\left(
\begin{array}
[c]{cccc}%
0 & \eta_{\mathcal{V}_{fib}\omega_{fib}} & 0 & \eta_{\mathcal{V}_{fib}%
\omega_{den}}\\
-\eta_{\mathcal{V}_{fib}\omega_{fib}}^{t} & 0 & -\eta_{\mathcal{V}_{fib}%
\omega_{den}}^{t} & 0\\
0 & \eta_{\mathcal{V}_{den}\omega_{fib}} & 0 & \eta_{\mathcal{V}_{den}%
\omega_{den}}\\
-\eta_{\mathcal{V}_{den}\omega_{fib}}^{t} & 0 & -\eta_{\mathcal{V}_{den}%
\omega_{den}}^{t} & 0
\end{array}
\right)
\end{equation}
where
\begin{equation}
\eta_{\mathcal{V}\omega}^{t}(\mathsf{z},\mathsf{z}^{\prime})=\eta
_{\omega\mathcal{V}}(\mathsf{z}^{\prime},\mathsf{z})
\end{equation}
The expectation values occuring in eq. $\left(  \ref{vomega}\right)  $ are
integral kernels of the unnormalized $N$-body density operator
\begin{align}
D^{N}\left(  \gamma\right)   &  =\left|  \Psi(\mathcal{V}_{fib}(\mathcal{V}%
_{den}),\omega_{fib}(\mathcal{V}_{den}),\mathcal{V}_{den},\omega\left(
\mathcal{V}_{den}\right)  \right\rangle \nonumber\\
&  \qquad\qquad\times\left\langle \Psi(\mathcal{V}_{fib}(\mathcal{V}%
_{den}),\omega_{fib}(\mathcal{V}_{den}),\mathcal{V}_{den},\omega\left(
\mathcal{V}_{den}\right)  \right|
\end{align}
and its associated reduced density operators. For instance
\begin{equation}
D^{N}\left(  \mathcal{T}\left(  \gamma\left(  t\right)  \right)  \right)
\left(  \mathsf{z},\mathsf{z}\right)  =\left(  -1\right)  ^{N_{e}+N_{n}%
}\left\langle \Psi\left(  \mathcal{T}\left(  \gamma\left(  t\right)  \right)
\right)  |\Lambda\left(  \mathsf{z}\right)  \Psi\left(  \mathcal{T}\left(
\gamma\left(  t\right)  \right)  \right)  \right\rangle
\end{equation}
when all the coordinates contained in $\mathsf{z}$ are distinct. If $N_{e}-p$
of the electronic coordinates are the same then one obtains the integral
kernel of areduced density operator referring to $p$ electrons, while if
$N_{n}-q$ of the nuclear coordinates are the identical for fermionic nuclei
one obtains the integral kernel of a reduced density operator referring to
$N_{n}-q$ nuclei. For bosonic nuclei one does not produce reduced density
operators as the order of the field operators is never reduced. The kernels
produced from the operators $\Lambda\left(  \mathsf{z}\right)  \Lambda\left(
\mathsf{z}^{\prime}\right)  $ are also of the $N$-body density operator and
its associated reduced density operators. In order to produce non zero kernels
for electronic and either fermionic nuclear or distinguishable nuclear systems
$\mathsf{z=z}^{\prime}$ (see Appendix A) and then $\Lambda\left(
\mathsf{z}\right)  ^{2}=\Lambda\left(  \mathsf{z}\right)  $.

One can systematically consdier the contributions to the metric $\eta$
arrising from the full and various reduced density operators by decomposing
$\mathcal{\breve{T}}$ into \textit{pure particle parts }$\left\{
\mathcal{\breve{T}}_{pq}\right\}  $, $p$-electron and $q$-nuclear, defined for
electrons and nuclei as%
\begin{equation}
\mathcal{\breve{T}}_{pq}=\int\mathcal{T}_{pq}\left(  \mathsf{x,X}\right)
\prod_{1\leq i\leq p}\Phi_{e}^{\dagger}(x_{i})\Phi_{e}(x_{i})\prod_{1\leq
j\leq q}\Phi_{n}^{\dagger}(X_{j})\Phi_{n}(X_{j})\prod_{1\leq i\leq p}%
dx_{i}\prod_{1\leq j\leq q}dx_{j}%
\end{equation}
where $\left\{  x_{1},\ldots,x_{p}\right\}  $ and $\left\{  X_{1},\ldots
,X_{q}\right\}  $ have distinct values over the integration range. The
operator $\mathcal{T}$ can be decomposed as
\begin{equation}
\mathcal{\breve{T}=}\sum_{0\leq p\leq N_{e};0\leq q\leq N_{n}}\mathcal{\breve
{T}}_{pq}%
\end{equation}
where $\mathcal{\breve{T}}_{0q}$ and $\mathcal{\breve{T}}_{p0}$ refer to
electronic and nuclear parts that are purely time dependent i.e. ones that
have no spatial nor spin dependence. Each of the pure parts $\mathcal{\breve
{T}}_{pq}$ can be expressed as
\begin{equation}
\mathcal{\breve{T}}_{pq}=\mathcal{\breve{V}}_{fib\,pq}+i\breve{\omega
}_{fib\,pq}+\mathcal{\breve{V}}_{den\,pq}+i\breve{\omega}_{den\,pq}%
\end{equation}
and the variables $\left\{  \mathcal{\breve{V}}_{fib},\breve{\omega}%
_{fib},\mathcal{\breve{V}}_{den},\breve{\omega}_{den}\right\}  $ can be varied
and approximated in terms of their independent pure parts
\begin{equation}
\left\{  \left\{  \mathcal{\breve{V}}_{fib\,pq}\right\}  ,\left\{
\breve{\omega}_{fib\,pq}\right\}  ,\left\{  \mathcal{\breve{V}}_{den\,pq}%
\right\}  ,\left\{  \breve{\omega}_{den\,pq}\right\}  ;0\leq p\leq N_{e},0\leq
q\leq N_{n}\right\}
\end{equation}
hence
\begin{equation}
e^{\mathcal{\breve{T}}}=\prod_{0\leq p\leq N_{e};0\leq q\leq N_{n}%
}e^{\mathcal{\breve{T}}_{pq}}%
\end{equation}
and we have analogous expressions for $\left\{  \mathcal{\breve{V}}%
_{fib},\breve{\omega}_{fib},\mathcal{\breve{V}}_{den},\breve{\omega}%
_{den}\right\}  $. These decompositions will be used as basic template for the
systematic approximations scheme for the construction of approximate time
dependent functionals that we are developing in this article.

The variations with respect to the different irreducible components involve
metric terms based on combinations of different reduced density operators. In
order to illuminate the details of these terms we concentrate on the
variations of the components $\left\{  0\leq p\leq2,0\leq q\leq2\right\}  $

\section{Computation of Metric and Force Terms}

\section{Summary Remarks}

\section{Acknowledgements}

We thank \ldots SBT was supported in part by the U.S. National Science Foundation.

\section{Appendix A Properties of Field Operators}

Fermionic local density operators are idempotent and commute with each other
i.e.
\begin{align}
&  \left(  \Phi_{\lambda}^{\dagger}(\kappa)\Phi_{\lambda}(\kappa)\right)
^{2}=\Phi_{\lambda}^{\dagger}(\kappa)\Phi_{\lambda}(\kappa)\nonumber\\
&  \left[  \Phi_{\mathcal{\lambda}}^{\dagger}(\kappa_{1})\Phi
_{\mathcal{\lambda}}(\kappa_{1}),\Phi_{\mathcal{\lambda}}^{\dagger}(\kappa
_{2})\Phi_{\mathcal{\lambda}}(\kappa_{2})\right]  =0
\end{align}
where $\kappa$ can be electron or nuclear coordinates associated with
$\lambda=e$ or $n$. Products of these local density operators have the
decomposition property can be rearranged into nromal form as
\begin{align}
&  \prod_{1\leq i\leq m}\Phi_{\mathcal{\lambda}}^{\dagger}(\kappa_{i}%
)\Phi_{\mathcal{\lambda}}(\kappa_{i})=\nonumber\\
&  =\sum_{1\leq i_{1}<i_{2}\leq m}\left\{  \Phi_{\mathcal{\lambda}}^{\dagger
}(\kappa_{i_{1}})\Phi_{\mathcal{\lambda}}(\kappa_{i_{2}})\prod
_{\substack{1\leq j_{1}<j_{2}\leq m \\j_{1},j_{2}\neq i_{1},i_{2} }}
\delta\left(  \kappa_{j_{1}}-\kappa_{j_{2}}\right)  \right\}  -\nonumber\\
&  \sum_{1\leq i_{1}<i_{2}<i_{3}<i_{4}\leq m}\left\{  \Phi_{\mathcal{\lambda}%
}^{\dagger}(\kappa_{i_{1}})\Phi_{\mathcal{\lambda}}^{\dagger}(\kappa_{i_{2}%
})\Phi_{\mathcal{\lambda}}(\kappa_{i_{2}})\Phi_{\mathcal{\lambda}}%
(\kappa_{i_{1}})\right. \nonumber\\
&  \qquad\qquad\qquad\qquad\times\left.  \prod_{\substack{1\leq j_{1}%
<j_{2}\leq m \\j_{1},j_{2}\neq i_{1},i_{2},i_{3},i_{4} }} \delta\left(
\kappa_{j_{1}}-\kappa_{j_{2}}\right)  \right\}  +\cdots\nonumber\\
&  \left(  -1\right)  ^{r}\sum_{1\leq i_{1}<\cdots<i_{r}\leq m}\left\{
\Phi_{\mathcal{\lambda}}^{\dagger}(\kappa_{i_{1}})\cdots\Phi_{\mathcal{\lambda
}}^{\dagger}(\kappa_{i_{r}})\Phi_{\mathcal{\lambda}}(\kappa_{i_{r}})\cdots
\Phi_{\mathcal{\lambda}}(\kappa_{i_{1}})\right. \nonumber\\
&  \qquad\qquad\qquad\times\left.  \prod_{\substack{1\leq j_{1}<j_{2}\leq m
\\j_{1},j_{2}\neq i_{1},\ldots,,i_{r} }} \delta\left(  \kappa_{j_{1}}%
-\kappa_{j_{2}}\right)  \right\}  +\nonumber\\
&  \qquad\qquad\cdots\left(  -1\right)  ^{m}\Phi_{\mathcal{\lambda}}^{\dagger
}(\kappa_{1})\cdots\Phi_{\mathcal{\lambda}}^{\dagger}(\kappa_{m}%
)\Phi_{\mathcal{\lambda}}(\kappa_{m})\cdots\Phi_{\mathcal{\lambda}}(\kappa
_{1}) \label{decomp}%
\end{align}
Thus one can observe that products of more than $N_{e}$ electron operators and
$N_{n}$ nuclear operators (belonging to the same fermionic species) on a
$N=N_{e}+N_{n}$ particle state can be duplicated by products of order $\leq$
$N_{e}$ and $\leq$ $N_{n}$. By noting that if coordinates are equal in the top
line of eq. $\left(  \ref{decomp}\right)  $ one can permute the product so
that local density operators with the same coordinates are adjacent to each
other and thus by the idempotency property produce just the first power of
that local density operator and hence see that such a product is a lower order
one containing distinct coordinates and thus conclude that
\begin{align}
\prod_{1\leq i\leq m}\Phi_{\mathcal{\lambda}}^{\dagger}(\kappa_{i}%
)\Phi_{\mathcal{\lambda}}(\kappa_{i})  &  =\prod_{1\leq i\leq m^{\prime}\leq
m}\Phi_{\mathcal{\lambda}}^{\dagger}(\kappa_{i})\Phi_{\mathcal{\lambda}%
}(\kappa_{i})\nonumber\\
&  =\left(  -1\right)  ^{m^{\prime}}\Phi_{\mathcal{\lambda}}^{\dagger}%
(\kappa_{1})\cdots\Phi_{\mathcal{\lambda}}^{\dagger}(\kappa_{m^{\prime}}%
)\Phi_{\mathcal{\lambda}}(\kappa_{m^{\prime}})\cdots\Phi_{\mathcal{\lambda}%
}(\kappa_{1})
\end{align}
Hence one only need to consider electronic and nuclear coordinates that are
distinct amongst the same species.

Products of Bosonic local density operators have a similiar decomposition to
eq. $\left(  \ref{decomp}\right)  $without alternating + and - signs i.e.
\begin{align}
&  \prod_{1\leq i\leq m}\Phi_{n}^{\dagger}(\kappa_{i})\Phi_{n}(\kappa_{i})=\\
&  \sum_{1\leq r\leq N_{n}}\sum_{1\leq i_{1}<\cdots<i_{r}\leq m}\left\{
\Phi_{\mathcal{\lambda}}^{\dagger}(\kappa_{i_{1}})\cdots\Phi_{\mathcal{\lambda
}}^{\dagger}(\kappa_{i_{r}})\Phi_{\mathcal{\lambda}}(\kappa_{i_{r}})\cdots
\Phi_{\mathcal{\lambda}}(\kappa_{i_{1}})\right. \\
&  \qquad\qquad\qquad\qquad\qquad\qquad\times\left.  \prod_{\substack{1\leq
j_{1}<j_{2}\leq m \\j_{1},j_{2}\neq i_{1},\ldots,,i_{r} }} \delta\left(
\kappa_{j_{1}}-\kappa_{j_{2}}\right)  \right\}
\end{align}
and the same observation applies to the boson case as the fermion one. These
operators still commute but are no longer idempotent i.e
\begin{align}
&  \left(  \Phi_{\lambda}^{\dagger}(\kappa)\Phi_{\lambda}(\kappa)\right)
^{2}\neq\Phi_{\lambda}^{\dagger}(\kappa)\Phi_{\lambda}(\kappa)\nonumber\\
&  \left[  \Phi_{\mathcal{\lambda}}^{\dagger}(\kappa_{1})\Phi
_{\mathcal{\lambda}}(\kappa_{1}),\Phi_{\mathcal{\lambda}}^{\dagger}(\kappa
_{2})\Phi_{\mathcal{\lambda}}(\kappa_{2})\right]  =0
\end{align}
Products of bosonic local density operators do not reduce in particle number
when coordinates are equal due to this lack of idempotency.

Products, $\left\{  \Phi_{\mathcal{\lambda}_{1}}^{\dagger}(\kappa_{1}%
)\Phi_{\mathcal{\lambda}_{1}}(\kappa_{1}),\ldots,\Phi_{r}^{\dagger}(\kappa
_{r})\Phi_{r}(\kappa_{r})\right\}  $, of local density operators for
distinguishable species commute and do not have to fulfill any commutation
relationships so
\begin{align}
&  \left[  \Phi_{\mathcal{\lambda}_{i}}^{\dagger}(\kappa_{i})\Phi
_{\mathcal{\lambda}_{i}}(\kappa_{i}),\Phi_{\mathcal{\lambda}_{j}}^{\dagger
}(\kappa_{j})\Phi_{\mathcal{\lambda}_{j}}(\kappa_{j})\right]  =0\nonumber\\
&  \prod_{1\leq i\leq m}\Phi_{\mathcal{\lambda}_{i}}^{\dagger}(\kappa_{i}%
)\Phi_{\mathcal{\lambda}_{i}}(\kappa_{i})=\Phi_{\mathcal{\lambda}_{1}%
}^{\dagger}(\kappa_{i_{1}})\cdots\Phi_{\mathcal{\lambda}_{m}}^{\dagger}%
(\kappa_{i_{m}})\Phi_{\mathcal{\lambda}_{m}}(\kappa_{i_{m}})\cdots
\Phi_{\mathcal{\lambda}_{1}}(\kappa_{i_{1}})
\end{align}
By using the preceding relationships one can easily obtain reodering and
decomposition expressions for mixed fermionic-bosonic-distinguishable systems.

\textsc{Part I references are next}

\begin{thebibliography}{99}                                                                                               %


\bibitem {Teil92}J.~Teilhaber, Phys. Rev. B \textbf{46}, 12990 (1992).

\bibitem {Stener95}M.~Stener, P.~Decleva, and A.~Lisini, J. Phys. B [At. Mol.
Opt. Phys.] \textbf{28}, 4973 (1995).

\bibitem {Saalmann96}U.~Saalmann and R.~Schmidt, Z. Phys. D \textbf{38}, 153 (1996)

\bibitem {Tong97}X.M.~Tong and S-I.~Chu, Phys. Rev. A \textbf{55}, 3406 (1997).

\bibitem {Serra97a}L.~Serra and A.~Rubio, Z. Physik D \textbf{40}, 262 (1997).

\bibitem {Ullrich97}C.A.~Ullrich, P.G.~Reinhard, and E.~Suraud, J. Phys.B [At.
Mol. Optical Phys.] \textbf{30}, 5043 (1997).

\bibitem {AG97}A.~G\"orling, Phys. Rev. B \textbf{55}, 2630 (1997).

\bibitem {Vignale97}G.~Vignale, C.A.~Ullrich, and S.~Conti, Phys. Rev. Lett.
\textbf{79}, 4878 (1997).

\bibitem {Calvayrac98}F.~Calvayrac, A.~Domps, P.G.~Reinhard, E.~Suraud, and
C.A.~Ullrich, Euro. Phys. J. \textbf{4}, 207-18 (1998).

\bibitem {Reinhard98}P.G.~Reinhard, E.~Suraud, and C.A.~Ullrich, Euro. Phys.
J. \textbf{1}, 303-11 (1998).

\bibitem {Ullrich98}C.A.~Ullrich, P.G.~Reinhard, and E.~Suraud, Phys. Rev. A
\textbf{57}, 1938-43 (1998).

\bibitem {Telnov98}D.A.~Telnov and S-I.~Chu, Phys. Rev. A \textbf{58}, 4749-56 (1998).

\bibitem {Tong98}X.-M.~Tong and S.-I.~Chu, Phys. Rev. A \textbf{57}, 452 (1998).

\bibitem {Campillo99}I.~Campillo, A.~Rubio, and J.M.~Pitarke, Phys. Rev.B
\textbf{59}, 12188-91 (1999).

\bibitem {Vasiliev99}I.~Vasiliev, S.~Ogut, and J.R.~Chelikowsky, Phys. Rev.
Lett. \textbf{82}, 1919-22 (1999).

\bibitem {Knospe99}O.~Knospe, J.~Jellinek, U.~Saalmann, and R.~Schmidt, Eur.
Phys. J. D \textbf{5}, 1 (1999) and references therein.

\bibitem {Hessler}P.~Hessler, J.~Park, and K.~Burke, Phys. Rev. Lett.
\textbf{82}, 378 (1999).

\bibitem {Gross96}E.K.U.~Gross, J.F.~Dobson, and M.~Petersilka, in
\textit{Density Functional Theory II - Topics in Current Chemistry 181}, R.F.
Nalewajski ed. (Springer, Berlin and Heidelberg, 1996) pp 81-172.

\bibitem {Casida95}M.~Casida in D.P. Chong ed., \textit{Recent Advances in
Density Functional Methods} (World Scientific, Singapore, 1995) pp 155-192.

\bibitem {Bartolotti81}L.J.~Bartolotti, Phys. Rev. A \textbf{24}, 1661 (1981).

\bibitem {DebGhosh82}B.M.~Deb and S.K.~Ghosh, J. Chem. Phys. \textbf{77}, 342 (1982).

\bibitem {RG84}E.~Runge and E.K.U.~Gross, Phys. Rev. Lett. \textbf{52}, 997 (1984).

\bibitem {HK}P.~Hohenberg and W.~Kohn, Phys. Rev. \textbf{136}, B864 (1964).

\bibitem {Xu85}B-X.~Xu and A.K.~Rajagopal, Phys. Rev. A \textbf{31}, 2682 (1985).

\bibitem {Dhara87}A.K.~Dhara and S.K.~Ghosh, Phys. Rev. A \textbf{35}, 442 (1987).

\bibitem {Levy79}M.~Levy, Proc. Natl. Acad. Sci. USA \textbf{76}, 6062-65 (1979).

\bibitem {LevyPerdew85}M.~Levy and J.P.~Perdew in R.M.~Dreizler and J.~da
Providencia eds., \textit{Density Functional Methods in Physics} (Plenum,
N.Y., 1985) pp 11-30.

\bibitem {Lieb}E.H.~Lieb, Internat. J. Quantum. Chem. \textbf{24}, 243 (1983).

\bibitem {KD86}H.~Kohl and R.M.~Dreizler, Phys. Rev. Lett. \textbf{56}, 1993 (1986).

\bibitem {Raj94A}A.K.~Rajagopal, Phys. Lett. A \textbf{193}, 363 (1994).

\bibitem {RajBuot94B}A.K.~Rajagopal and F.A.~Buot, Phys. Rev. E \textbf{50},
721 (1994).

\bibitem {LiTong86}T-C.~Li and P-q. Tong, Phys. Rev. A \textbf{34}, 529 (1986).

\bibitem {Qian98}A.~Qian and V.~Sahni, Phys. Lett. A \textbf{247}, 303 (1998).

\bibitem {Chermyak00}V.~Chernyak and S.~Mukamel J. Chem. Phys. \textbf{112},
3572 (2000)

\bibitem {DeumensA}E.~Deumens, A.~Diz, R.~Longo, and Y.~{\"O}hrn, Rev. Mod.
Phys. \textbf{66}, 917 (1994).

\bibitem {DeumensB}E.~Deumens and Y.~{\"O}hrn, J. Chem. Soc., Faraday Trans.
\textbf{93}, 919 (1997) and refs.~therein.

\bibitem {BLWSBT99a}B.~Weiner and S.B.~Trickey, Adv. Quantum Chemistry
\textbf{35} (Academic, San Diego, 1999). 217.

\bibitem {BLWSBT00a}``State Energy Functionals and Variational Equations in
Density Functional Theory,'' B.~Weiner and S.B.~Trickey, J. Molec. Struct.
(Theochem) \textbf{501}-\textbf{02}, 65 (2000).

\bibitem {Leeuwen99}R.~van Leeuwen, Phys. Rev. Lett. \textbf{82}, 3863 (1999).

\bibitem {GrossKohn90}W.~Kohn and E.K.U.~Gross in \textit{Density Functional
Theory for Many-Fermion Systems} S.B. Trickey spec. ed., Adv. Quant. Chem.
\textbf{21} (Academic, San Diego, 1990), 255.

\bibitem {Fois88}E.S.~Fois, A.~Selloni, M.~Parinello,and R.~Car, J. Phys.
Chem. \textbf{92}, 3268 (1988)

\bibitem {PerdewSavin}J.P.~Perdew, A.~Savin, and K.~Burke, Phys. Rev. A
\textbf{51}, 4531 (1995).

\bibitem {BLWSBT98}B.~Weiner and S.B.~Trickey, Internat. J. Quantum Chem.
\textbf{69}, 451 (1998).

\bibitem {SavinSem2}A.~Savin, in J.M. Seminario ed., \textit{Recent
Developments and Applications of Modern Density Functional Theory} (Elsevier,
Amsterdam, 1996), p. 351.

\bibitem {Sjoqvist}E.~Sj\"oqvist and M.~Hedstr\"om, Phys. Rev. A \textbf{56}
3417 (1997).

\bibitem {Bruhat}\textit{Analysis, Manifolds, and Physics}, Y.~Choquet-Bruhat
and C. DeWitt-Morette with M. Dillrd-Bleick, Revised Edition (North- Holland,
Amsterdam, 1982) p.350 ff.

\bibitem {affinebundleref}\textit{Encyclopedic Dictionary of Mathematics},
K.~It\^o ed. (MIT Press, Cambridge, 1996) vol. 1, p. 569

\bibitem {Harriman81}J. Harriman, Phys. Rev. A \textbf{24}, 680 (1981).

\bibitem {KramSar}P.~Kramer and M.~Saraceno, \textit{Geometry of the Time
Dependent Variational Principle in Quantum Mechanics} (Springer, New York 1981).

\bibitem {etainv}\textit{Classical Dynamics A Modern Perspective},
E.C.G.~Sudarshan and N.~Mukunda (John Wiley and Sons, NY, 1974) 112-19.

\bibitem {BLW94}B.~Weiner, E.~Deumens, and Y.~\"Ohrn, J. Math. Phys.
\textbf{35}, 1139 (1994).

\bibitem {Ortiz81}J.V.~Ortiz, B.~Weiner, and Y.~\"Ohrn, Internat J. Quantum
Chem. S-\textbf{15}, 113 (1981).

%\bibitem{coneref} {\sc Need Reference here}


\bibitem {Kummer}H.~Kummer, J. Math. Phys. \textbf{8}, 2063 (1967)

\bibitem {BerezinShubin}F.A.~Berezin and M.A..~Shubin, \textit{The
Schr\"odinger Equation} (Kluwer Academic Press, ??? 1981).

\bibitem {Richtmyer}R.~Richtmyer, \textit{Principles of Advanced Mathematical
Physics} (Springer Verlag, NY, 1978) vol. 1, p. 245.
\end{thebibliography}


\end{document}